\section{Testing}

\subsection{Testing Methodology}
The repository contains extensive diagnostic logging and clearly defined RTOS
periods, but it does not contain a full benchmark harness, oscilloscope traces,
or deadline-monitor instrumentation. Therefore, the validation in this report is
based on two evidence sources:
\begin{enumerate}
  \item \textbf{code-driven timing analysis}, derived from task periods, queue
  usage, and blocking behavior; and
  \item \textbf{serial-log validation}, derived from the diagnostic messages
  emitted by the firmware such as \texttt{[Sensor]}, \texttt{[LCD]},
  \texttt{[Web]}, \texttt{[TinyML]}, and \texttt{[CoreIoT]}.
\end{enumerate}

This means the project can be evaluated confidently as a \emph{soft real-time}
embedded system, but not yet as a formally verified hard real-time system.

\subsection{Test Plan}
The implemented test plan covers all assignment tasks and the interactions
between them:
\begin{itemize}
  \item verify periodic sensor acquisition and DHT20 initialization;
  \item verify temperature-to-LED blink mapping;
  \item verify humidity-to-NeoPixel color behavior;
  \item verify LCD rendering under normal, AP, STA, and TinyML states;
  \item verify AP dashboard response and external LED control;
  \item verify Wi-Fi scan, profile save, STA connect, timeout fallback, and
  BOOT-button AP recovery;
  \item verify TinyML batching and prediction publication;
  \item verify CoreIOT connection, RPC subscription, and telemetry publishing.
\end{itemize}

\subsection{Representative Log Evidence}
The following diagnostic messages are important for system-level testing:
\begin{itemize}
  \item \texttt{[Sensor] DHT20 initialized.}
  \item \texttt{[Sensor] T=\ldots, H=\ldots}
  \item \texttt{[LCD] LCD initialized.}
  \item \texttt{[Web] AP mode started \ldots}
  \item \texttt{[Web] STA connect start \ldots}
  \item \texttt{[Web] STA connected \ldots}
  \item \texttt{[Web] STA connect timeout -> back to AP}
  \item TinyML logs containing averaged values and rain probability
  \item \texttt{[CoreIoT] Attempting MQTT connection\ldots}
  \item \texttt{[CoreIoT] Subscribe v1/devices/me/rpc/request/+ -> OK}
  \item \texttt{[CoreIoT] Publish -> OK}
\end{itemize}

Even without an external benchmark framework, these logs give a clear
observation path for validating control-flow transitions and subsystem
activation.

\subsection{Timing Analysis by Task}
Table~\ref{tab:timing-analysis} summarizes the expected timing behavior of the
main tasks based on the implemented code.

\begin{table}[!htbp]
\centering
\caption{Timing-oriented analysis of the implemented tasks}
\label{tab:timing-analysis}
\begin{tabular}{p{3cm}p{3.2cm}p{7.8cm}}
\toprule
\textbf{Task} & \textbf{Expected timing} & \textbf{Observation and interpretation} \\
\midrule
Sensor task & 2 s fixed period & Uses \texttt{vTaskDelayUntil()}, which is the strongest timing primitive in the project. This should keep long-term drift low as long as execution time stays below the period. \\
LED task & event-driven, mode-dependent half-period & Responds immediately to queue input because \texttt{xQueueReceive()} wakes as soon as a command arrives. Blink cadence depends on the current mode. \\
NeoPixel task & event-driven & Updates on each humidity command. Worst-case refresh is bounded by the sensor period, because humidity changes originate from the sensor task. \\
LCD task & event-driven, queue wait up to 250 ms & Typical update occurs soon after a new sample arrives. Additional updates occur on Wi-Fi RSSI or display-state changes. \\
Web task & 20 ms service loop & Appropriate for responsive local HTTP handling. Human interaction latency should feel near-instant on a local AP. \\
TinyML task & one inference per 10 samples & Since sensing is 2 s periodic, inference is generated approximately every 20 s. This is the main intentional latency in the system. \\
CoreIOT task & 10 s publish period & Telemetry is sent every 10 s after Wi-Fi and MQTT readiness. Network latency is nondeterministic and outside RTOS control. \\
\bottomrule
\end{tabular}
\end{table}

\subsection{Evaluation of RTOS Timing Quality}
From a real-time perspective, the strongest part of the implementation is the
sensing pipeline. The sensor task is explicitly periodic, and all downstream
tasks are primarily event-driven. This is a good architecture because the
periodic source is stable while the rest of the system reacts to the arrival of
new data.

The following observations can be made:
\begin{itemize}
  \item The project does \emph{not} rely on a long blocking loop that mixes all
  concerns together.
  \item The sensor task has the clearest timing contract in the whole system.
  \item The web task remains responsive because it executes at a short
  20 ms period.
  \item The cloud task is decoupled from sensing and therefore should not delay
  the DHT20 pipeline.
  \item The LCD and TinyML tasks trade immediacy for better stateful behavior.
\end{itemize}

However, it is also important to state what is \emph{not} yet proven:
\begin{itemize}
  \item Worst-case execution time (WCET) has not been measured for every task.
  \item There is no deadline-monitor task or explicit miss counter.
  \item Cloud and Wi-Fi latencies remain network-dependent.
  \item TinyML inference time is not timestamped in the logs as a numeric
  micro-benchmark.
\end{itemize}

Therefore, the system can be judged as meeting the assignment's practical RTOS
expectation, but it should be classified as a soft real-time implementation
rather than a formally validated hard real-time design.

\subsection{Task-by-Task Functional Verification}
\subsubsection*{Task 1: LED behavior.}
The LED task meets the requirement of at least three temperature-dependent
behaviors. The classification function distinguishes cold, normal, and hot
regions, while the manager task converts each region into a visibly different
blink rate. In addition, the module already includes event pulses, which
improves debuggability and future extensibility.

\subsubsection*{Task 2: NeoPixel behavior.}
The NeoPixel task satisfies the humidity mapping requirement and goes beyond it
by providing a continuous gradient rather than only three static colors. Because
the task receives data through a queue, synchronization with the sensor task is
clean and non-blocking.

\subsubsection*{Task 3: LCD monitoring.}
The LCD subsystem clearly satisfies the requirement of multiple display states.
It supports sensor error feedback, AP/STA context awareness, and TinyML
forecast display. The use of an I\textsuperscript{2}C mutex is technically
important and directly aligned with the synchronization emphasis of the
assignment.

\subsubsection*{Task 4: Web server.}
The web task strongly satisfies the assignment. It provides:
\begin{itemize}
  \item Wi-Fi scan,
  \item Wi-Fi save and recall,
  \item connect and fallback flow,
  \item an external LED on/off button,
  \item an external LED brightness slider,
  \item live sensor data display, and
  \item a recent-history chart.
\end{itemize}
The implementation therefore exceeds the minimum requirement of ``two buttons''
by including a meaningful monitoring dashboard and configuration workflow.

\subsubsection*{Task 5: TinyML.}
The TinyML module fulfills the requirement of deploying a model on the device
and exposing the result to the rest of the system. The current evaluation is
primarily architectural and functional rather than dataset-statistical, because
the repository focuses on deployment code. Still, the task is fully integrated
and produces live inference output in the running firmware.

\subsubsection*{Task 6: CoreIOT.}
The CoreIOT task satisfies the central cloud requirement: publishing occurs in
STA mode after Internet readiness is signaled. The task also subscribes to RPC
messages, which is a strong implementation detail because it turns cloud
communication into a bidirectional control channel.

\subsection{End-to-End Responsiveness}
Looking at the end-to-end flow, the expected delays are as follows:
\begin{itemize}
  \item \textbf{Sensor to LED/NeoPixel:} effectively one sensor period plus
  immediate queue wake-up, so around the 2 s sensing cadence.
  \item \textbf{Sensor to LCD:} typically within the same sampling cycle, plus
  a small scheduling/rendering delay.
  \item \textbf{Sensor to Web dashboard:} one sensor period plus the browser's
  2 s polling interval.
  \item \textbf{Sensor to TinyML result:} approximately 20 s because of
  batching.
  \item \textbf{Sensor to CoreIOT publish:} bounded by the 10 s publish period,
  assuming Wi-Fi and MQTT are already connected.
\end{itemize}

These values are consistent with the intended operating scope of the system.
The implementation is optimized for functional integration and predictable
behavior rather than ultra-low-latency inference.

\subsection{Current Limitations of the Test Evidence}
The current code base supports convincing functional testing, but a more
advanced evaluation could still improve the report:
\begin{itemize}
  \item add explicit timestamped logs for task start/end times,
  \item measure TinyML inference duration directly,
  \item collect repeated Wi-Fi reconnect statistics,
  \item record actual CoreIOT publish success rates over long runs,
  \item export charts and screenshots from the AP dashboard.
\end{itemize}

Even without these additions, the current implementation provides convincing
evidence that all major assignment tasks have been achieved and that the RTOS
design behaves consistently with its intended timing model.
